% Cover letter (PDF upload) -- The Visual Computer (TVCJ).
% Standalone; lmodern+T1 to avoid pk-bitmap generation in this sandbox.
% Content mirrors cover_letter_tvcj.txt verbatim.
\documentclass[10pt]{article}
\usepackage{lmodern}
\usepackage[T1]{fontenc}
\usepackage[a4paper,margin=0.85in]{geometry}
\setlength{\parindent}{0pt}
\setlength{\parskip}{0.55em}
\setlength{\emergencystretch}{2em}
\usepackage[hidelinks]{hyperref}
\hypersetup{
  pdfauthor={Hainan Huang},
  pdftitle={Cover Letter - The Visual Computer}
}
\pagestyle{empty}

\begin{document}

\begin{flushright}
September 29, 2026
\end{flushright}

\noindent Dear Editor-in-Chief and Editors of The Visual Computer,

\medskip
\noindent We are pleased to submit our manuscript entitled ``Dynamic Prototype Routing for
Lightweight Image Super-Resolution'' for consideration as a research article in
\emph{The Visual Computer}.

The journal's current scope explicitly includes ``Image and Video Processing,
Restoration, and Enhancement'' at the intersection of computer vision, artificial
intelligence, and graphics. Single image super-resolution directly addresses this area by
recovering fine visual detail for display and downstream analysis. The journal has also
continued to publish lightweight and efficient single-image super-resolution studies in
recent volumes. Our manuscript addresses the same problem setting while contributing a
distinct mechanism for confidence-aware token routing.

CATANet's Token Aggregation Block groups related tokens before attention but depends on
cross-batch EMA centers, hard assignment labels, and a coupled prototype
update/confirmation step. We replace this router with Dynamic Prototype Routing (DPR), a
drop-in module that preserves the backbone and external tensor interface. DPR forms
per-image prototypes from current tokens, uses a learnable query bank to decouple prototype
generation from confirmation, retains assignment confidence in token ordering, global
gating, and a soft fallback, and combines a bounded routing temperature with a
usage-balancing regularizer to discourage prototype-slot collapse.

Across five benchmarks at x2/x3/x4, DPRNet matches CATANet-level accuracy with small gains
on most benchmark cells; at x4, it obtains 26.89~dB on Urban100. DPR adds a modest cost over
CATANet (660K vs 535K parameters; 52.14G vs 46.8G Multi-Adds), so we do not claim an
efficiency advantage over CATANet. It remains in the lightweight regime and uses fewer
Multi-Adds than SwinIR-light and SRFormer-light under the same input convention. Three-seed
ablations show that dynamic prototypes outperform a matched EMA history-center variant and
that the balancing loss improves both the balance and cross-seed stability of prototype
usage. Routing diagnostics further expose the learned temperature, confidence distribution,
usage entropy, and per-block cluster assignments.

Because DPR is a plug-and-play, interface-preserving module, it is directly reusable in
other window-attention visual-computing backbones. We believe this combination---a reusable
content-routing component, an honest efficiency/accuracy characterization, and
mechanism-level, visually interpretable analysis for image restoration---fits the stated
scope and readership of The Visual Computer.

This manuscript is original, has not been published previously, and is not under
consideration for publication elsewhere. All authors have read and approved the manuscript
and agree to its submission. The authors declare no competing interests and no external
funding. All datasets used in the study are publicly available; no new data were generated.

We suggest the following areas of expertise for reviewers: lightweight/efficient image
super-resolution, vision Transformers for image restoration, and content-aware token
routing / clustering.

Thank you for your time and consideration. We look forward to your response.

\medskip
\noindent Sincerely,

\medskip
\noindent Hainan Huang (corresponding author), on behalf of all authors \\
School of Computer Science, China University of Geosciences (Wuhan) \\
Future City Campus, No.\ 68 Jincheng Street, East Lake High-tech Development Zone, \\
Wuhan, Hubei 430078, China \\
\texttt{huanghainan.7@cug.edu.cn}

\end{document}
